% ------------------------------------------------------------
% A6.x Substrate Uplift Module B — TARTAN Projection Algebra (v0.1)
% Author: Borealis S. Hedling
% Altitude: A6.x (Substrate Uplift Layer)
% Purpose: Define the TARTAN projection algebra required for
%          modal adjacency, spectral continuity, and collapse
%          compatibility across the A6 manifold.
% ------------------------------------------------------------

\documentclass[11pt]{article}
\usepackage{amsmath, amssymb, amsthm}

\title{A6.x Substrate Uplift Module B: TARTAN Projection Algebra}
\author{Borealis S. Hedling}
\date{October 2026}

\begin{document}
\maketitle

\section*{Overview}

This module defines the TARTAN projection algebra required for A6 altitude
geometry. TARTAN provides the modal adjacency, projection continuity, and
collapse-surface compatibility required by the A6 Constructive Fixes and the
A7 Composite Substrate Operator Algebra.

% ------------------------------------------------------------
\section{TARTAN Modal Projection}

Let $\mu : M \to \mathcal{M}$ be the modal map.

\textbf{Definition (TARTAN Projection).}
\[
\Pi_{\mathrm{T}}(p) = \mu^{-1}(\mu(p)).
\]

This defines the modal fiber through $p$.

\textbf{Theorem (Adjacency Preservation).}
If $p^+ \sim p^-$, then:
\[
\mu(p^+) \in \mathcal{N}(\mu(p^-)).
\]

This enforces Fix 4 (modal adjacency).

% ------------------------------------------------------------
\section{TARTAN Spectral Projection}

Define the spectral operator:
\[
L_{\mathrm{T}} : TM \to TM.
\]

\textbf{Definition (Spectral Projection).}
\[
\Pi_{\mathrm{spec}}(p) = \mathrm{Spec}(L_{\mathrm{T}}(p)).
\]

\textbf{Theorem (Spectral Continuity).}
\[
\mathrm{Spec}(L_{\mathrm{T}}(p)) \text{ continuous in } p.
\]

This enforces Fix 5 (spectral continuity).

% ------------------------------------------------------------
\section{Spectral Mismatch Bound}

Define spectral mismatch:
\[
\delta_{\mathrm{spec}}(p) = \|\Pi_{\mathrm{spec}}(p) - \Pi_{\mathrm{spec}}(q)\|.
\]

\textbf{Constraint.}
\[
\delta_{\mathrm{spec}}(p) < \epsilon.
\]

This enforces Fix 6 (spectral mismatch bound).

% ------------------------------------------------------------
\section{TARTAN Collapse Compatibility}

Let $\Sigma$ be the collapse surface.

\textbf{Definition (TARTAN Collapse Projection).}
\[
\Pi_{\Sigma}(p) = 
\begin{cases}
p, & p \in \Sigma \\
\text{nearest point in } \Sigma, & p \notin \Sigma.
\end{cases}
\]

\textbf{Theorem (Local Detectability).}
\[
p \in \Sigma \Rightarrow \exists \epsilon > 0 :
B_\epsilon(p) \cap \Sigma \neq \emptyset.
\]

This enforces Fix 8 (collapse locality).

% ------------------------------------------------------------
\section{TARTAN Stability Envelope}

Define:
\[
E_{\mathrm{T}}(p) = \|\Pi_{\mathrm{T}}(p)\|.
\]

\textbf{Theorem (Monotonicity).}
If $E(p) \le E(q)$, then stability lifts monotonically.

This enforces Fix 9 (stability envelope monotonicity).

% ------------------------------------------------------------
\section{TARTAN Modal Lift}

Modal truth must lift through TARTAN:
\[
\mu(p) \leadsto \mu(q).
\]

\textbf{Theorem (Lift Preservation).}
If $\Pi_{\mathrm{T}}(p) = \Pi_{\mathrm{T}}(q)$, then modal truth lifts.

This enforces Fix 10 (modal lift preservation).

% ------------------------------------------------------------
\section{Conclusion}

The TARTAN projection algebra defined here provides the modal adjacency,
spectral continuity, collapse compatibility, and stability structures required
for A6 Constructive Fixes and A7 Composite Substrate Operator Algebra. It is
the second foundational substrate module in the A6.x uplift sequence.

\end{document}
